% Generated draft: requires mathematical and source review.
\documentclass[11pt]{article}
\usepackage[margin=25mm]{geometry}
\usepackage{fontspec}
\setmainfont{Latin Modern Roman}
\usepackage{amsmath,amssymb,mathtools,unicode-math}
\setmathfont{Latin Modern Math}
\usepackage{enumitem,needspace,xurl,hyperref}
\setlength{\parindent}{0pt}
\setlength{\parskip}{5pt}
\setlength{\emergencystretch}{3em}
\newcommand{\N}{\mathbb{N}}
\newcommand{\Z}{\mathbb{Z}}
\newcommand{\Q}{\mathbb{Q}}
\newcommand{\R}{\mathbb{R}}
\newcommand{\C}{\mathbb{C}}
\newcommand{\F}{\mathbb{F}}
\newcommand{\A}{\mathbb{A}}
\newcommand{\PP}{\mathbb{P}}
\newcommand{\E}{\mathbb{E}}
\newcommand{\eps}{\varepsilon}
\newcommand{\dd}{\,\mathrm{d}}
\begin{document}
\section*{UnsolvedMath ID: 701362}
\textbf{Source problem code:} catalogue entry
\subsection{Definitions and mathematical statement}
Let \(\mathbb P_{\mathrm{odd}}\) be the set of odd primes. For each subset \(\mathcal P\subseteq\mathbb P_{\mathrm{odd}}\), define the permitted positive moduli by
\[
M_{\mathcal P}=\{m\in\Z_{>0}:\text{every prime divisor of }m\text{ belongs to }\mathcal P\cup\{2\}\}.
\]
Following the supplied source terminology, let \(T_{\mathcal P}\) consist of all class transpositions in \(\mathrm{CT}(\Z)\) that interchange residue classes whose moduli both belong to \(M_{\mathcal P}\), and put
\[
\mathrm{CT}_{\mathcal P}(\Z)=\langle T_{\mathcal P}\rangle\leq\mathrm{CT}(\Z).
\]
The supplied record does not define the ambient group \(\mathrm{CT}(\Z)\) or the precise permutations called class transpositions. These terms therefore retain their intended source meanings; the record alone does not provide a self-contained construction of these groups.

The question is whether
\[
\forall\mathcal P,\mathcal Q\subseteq\mathbb P_{\mathrm{odd}},\qquad
\mathrm{CT}_{\mathcal P}(\Z)\cong\mathrm{CT}_{\mathcal Q}(\Z)
\ \Longrightarrow\ \mathcal P=\mathcal Q,
\]
where \(\cong\) denotes isomorphism of abstract groups. No finiteness or nonemptiness restriction on the prime sets is stated. The source reports that every \(\mathrm{CT}_{\mathcal P}(\Z)\) is simple, citing the publication listed below; this is supplied background, not an independently verified assertion.
\subsection{Short English statement}
S. Kohl asks whether distinct sets of odd primes always yield nonisomorphic groups generated by the permitted class transpositions of \(\Z\). The supplied registry labels the problem open. Precise definitions of the ambient group and class transpositions must be recovered from the source before this draft is self-contained.
\subsection{Sources}
\begin{itemize}
\item Local UnsolvedMath registry: supplied complete source record, permanent ID \(701362\), exact source code \texttt{catalogue entry}, title ``Kourovka Notebook Problem 17.60.'' The supplied status is \texttt{open}; this label has not been independently verified.
\item S. Kohl, proposer, 2010. The Kourovka Notebook (Unsolved Problems in Group Theory), Problem 17.60, 17th issue (2010); the record also identifies the 21st issue edition (2026), \texttt{arXiv:1401.0300v48}. The supplied collection description names E. I. Khukhro and V. D. Mazurov as editors.
\item Supplied reference for the simplicity assertion: \emph{Math. Z.}, volume \(264\), no. \(4\) (2010), pp.\ 927--938. The record supplies neither an article title nor author names for this reference. No source was independently consulted.
\end{itemize}
\end{document}
